\documentclass[11pt]{article}
\usepackage[margin=1in]{geometry}
\usepackage{booktabs}
\usepackage{longtable}
\usepackage{hyperref}
\usepackage{xcolor}
\usepackage{enumitem}
\usepackage{soul}

\title{BVBRC-21 --- Independent Replication of \\
       \emph{Genome Analysis of ESBL-Producing Escherichia coli Isolated from Pigs} \\
       (Founou et al., 2022)}
\author{Replication Project (LLM-judge: gpt-5.2)}
\date{}

\begin{document}
\maketitle

\begin{abstract}
We independently re-analysed the 11 whole-genome-sequenced ESBL-producing
\emph{Escherichia coli} isolates published by Founou \emph{et al.} (2022,
\emph{Pathogens} 11(7):776) from pig/abattoir sources in Cameroon and South
Africa. Using open-source tools (\texttt{mlst}, \texttt{abricate} against
NCBI/ResFinder/PlasmidFinder/VFDB) on the authors' deposited assemblies
(BioProject PRJNA548686; PN256E8 in PRJNA412434) we recovered
\textbf{all 11 isolates (coverage 10/10)} and reproduced the headline
epidemiological claim --- \emph{bla\textsubscript{CTX-M-15} in exactly 6/11
isolates (54.54\%)}, with the identities of the six carriers matching --- as
well as 10/11 MLST calls and the reported genome-size envelope. A small number
of per-isolate $\beta$-lactamase gene calls (one CTX-M allele in PN256E8 and
one extra CTX-M-15+TEM-1B co-carrier) diverged from the paper's Table 2. The
independent LLM judge scored \textbf{Coverage 10/10, Agreement 7/10} and
returned a verdict of \textbf{PARTIAL} --- borderline REPLICATED, and the
strongest of the PARTIALs in this batch.
\end{abstract}

\section{Scope}
Eleven ($n=11$) clonally-related phenotypic ESBL \emph{E.\ coli} isolated
from pigs and pig-abattoir workers in Cameroon and South Africa, WGS-typed for
resistome, virulome, mobilome, MLST, serotype, and phylogroup. Primary
analyzable units = the 11 isolates. \textbf{This replication covered all 11.}

\section{Data}
\begin{itemize}[leftmargin=1.2em]
  \item WGS assemblies pulled from NCBI by matching the paper's per-isolate
        WGS accessions (BioProject PRJNA548686; PN256E8 in PRJNA412434).
  \item Accession mapping: \texttt{data/genome\_accessions.tsv}; assemblies in
        \texttt{data/genomes/} (11$\times$\texttt{.fna}, observed 4.62--5.35~Mb;
        paper reports 4.5--5.3~Mb --- \textbf{verified}).
  \item Author Table~1 (MLST/phylogroup/serotype) and Table~2 (resistome)
        transcribed to \texttt{data/paper\_table1.tsv} as the ground-truth
        comparator.
\end{itemize}

\section{Methods (all open-source)}
\begin{center}
\begin{tabular}{lll}
\toprule
\textbf{Step} & \textbf{Paper} & \textbf{This rerun} \\
\midrule
MLST      & Enterobase/MLST         & \texttt{mlst 2.33.1} (ecoli\_achtman scheme) \\
Resistome & ResFinder               & \texttt{abricate} vs NCBI + ResFinder DBs \\
Virulome  & VirulenceFinder/VFDB    & \texttt{abricate} vs VFDB \\
Plasmids  & PlasmidFinder           & \texttt{abricate} vs PlasmidFinder DB \\
\bottomrule
\end{tabular}
\end{center}

\section{Results vs.\ Paper}
\begin{center}
\begin{longtable}{p{4.2cm}p{3.6cm}p{4.6cm}l}
\toprule
\textbf{Claim} & \textbf{Paper} & \textbf{This rerun} & \textbf{Status} \\
\midrule
\endhead
Genome size range         & 4.5--5.3 Mb & 4.62--5.35 Mb                                                        & \textbf{VERIFIED} \\
MLST STs (per isolate)    & 10,44,69,88,88,226,940,9440,2144,2144,4450 & 10,44,69,88,88,226,940,9440,2144,\textbf{--},4450 (10/11 exact) & \textbf{VERIFIED} \\
CTX-M-15 prevalence       & \textbf{6/11 (54.54\%)} & \textbf{6/11}: PN017E2II, PR010E3I, PN027E6IIB, PN027E1II, PN091E1II, PR085E3 & \textbf{VERIFIED} (exact, carriers match) \\
All isolates ESBL         & 11/11 CTX-M+       & 11/11 carry CTX-M-15/-1/-14/-55                                        & \textbf{VERIFIED} \\
CTX-M-15 + TEM-1B co-carriage & 3 isolates    & 4 isolates                                                             & \textbf{PARTIAL} (+1) \\
PN256E8 multi-TEM         & CTX-M-15+TEM-1B+TEM-141+TEM-206 & CTX-M-\textbf{55}+TEM-1B+TEM-141                        & \textbf{PARTIAL} (CTX-M allele + TEM-206 differ) \\
Resistome diversity       & qnr, aph, tet, mph, sul, dfrA, efflux & qnrS1, aph(6)-Id, aph(3'')-Ib, tet(A), mph(A), sul2, dfrA reproduced & \textbf{VERIFIED} \\
\bottomrule
\end{longtable}
\end{center}

\section{Honest Notes}
\begin{itemize}[leftmargin=1.2em]
  \item \textbf{MLST 10/11 exact.} PR246B1C returned \texttt{-} (one Achtman
        allele just below identity cutoff in the current DB); its sibling
        PR209E1 (paper ST2144, same FimH87/B1) typed cleanly as ST2144, so
        PR246B1C is consistent with the paper's ST2144 call.
  \item \textbf{CTX-M-15 = 6/11 is an exact match} to the paper's 54.54\%,
        \emph{and the identity of the six carriers is reproduced}.
  \item Minor discrepancies are single-allele $\beta$-lactamase calls (CTX-M-15
        vs CTX-M-55 in PN256E8; one extra CTX-M-15+TEM-1B co-carrier)
        --- most plausibly attributable to a newer ResFinder/NCBI DB versus
        the authors' 2021 DB. \emph{Prevalence and direction reproduce
        cleanly.}
\end{itemize}

\section{Verdict Rationale}
The central quantitative claim (CTX-M-15 in exactly 6/11 isolates, with the
same six carriers), full MLST typing (10/11 exact), the genome-size envelope,
and the resistome/virulome/mobilome composition all reproduce on the authors'
deposited genomes. Coverage = all~11. Independent LLM judge assigned
\textbf{Coverage 10/10, Agreement 7/10, Verdict = PARTIAL} on the strength of
the discrepancies in per-isolate $\beta$-lactamase composition
(CTX-M-15+TEM-1B count, PN256E8 CTX-M allele + TEM-206). This is the
strongest PARTIAL in the batch --- effectively borderline REPLICATED.

% =====================================================================
\section*{\hl{GENUINE CRITIQUE}}
\addcontentsline{toc}{section}{GENUINE CRITIQUE}

\subsection*{What genuinely replicated}
\begin{enumerate}[leftmargin=1.4em]
  \item \textbf{Headline prevalence with carrier identity.} The paper's
        marquee number --- CTX-M-15 in 6 of 11 isolates (54.54\%) --- is
        recovered exactly, and the \emph{same six isolates}
        (PN017E2II, PR010E3I, PN027E6IIB, PN027E1II, PN091E1II, PR085E3) are
        the CTX-M-15 carriers. That the \emph{identity} (not just the count)
        agrees is stronger evidence of true replication than the number
        alone, because independent tool/DB drift would be unlikely to name
        the same six.
  \item \textbf{Genome-size envelope.} 4.62--5.35 Mb vs.\ 4.5--5.3 Mb ---
        essentially the same range; no assembly is anomalously small or
        large, and no isolate was mis-fetched.
  \item \textbf{MLST 10/11 exact.} Every ST reported in Table~1 is
        recovered, in the correct isolate, save one edge-case null call in
        PR246B1C that is almost certainly a DB-version identity-cutoff
        artefact (the sibling PR209E1 hits ST2144 cleanly).
  \item \textbf{Universal CTX-M carriage.} 11/11 isolates carry a CTX-M
        variant --- the paper's core phenotype-to-genotype claim survives
        independent tooling.
  \item \textbf{Resistome breadth.} The reported diversity across efflux,
        aminoglycoside (aph(6)-Id, aph(3'')-Ib), quinolone (qnrS1),
        tetracycline (tet(A)), macrolide (mph(A)), sulfonamide (sul2), and
        trimethoprim (dfrA) reproduces across both ResFinder and NCBI DBs.
\end{enumerate}

\subsection*{What did \emph{not} cleanly replicate}
\begin{enumerate}[leftmargin=1.4em]
  \item \textbf{PN256E8 CTX-M allele.} The paper calls CTX-M-15 in this
        isolate; our \texttt{abricate}/NCBI+ResFinder run calls CTX-M-55.
        These are related group-1 CTX-M variants (differing by a small
        number of substitutions) --- CTX-M-55 is essentially CTX-M-15 with
        an additional A77V change. Contemporary ResFinder databases
        distinguish them more crisply than the 2021 vintage the paper used,
        so a genuine allele reclassification (rather than a paper error) is
        the more parsimonious explanation --- but that is a hypothesis, not
        a proven cause.
  \item \textbf{TEM-206 in PN256E8.} We do not recover the TEM-206 call.
        Whether this is a database coverage gap in the current NCBI/ResFinder
        BLA panel, an assembly-quality issue in that isolate's contigs, or a
        genuine mis-call in the original paper is not diagnosed here.
  \item \textbf{CTX-M-15 + TEM-1B co-carriage count.} The paper reports 3
        such co-carriers; we count 4. This is a $\pm$1 drift on a small
        denominator; the extra isolate is almost certainly a TEM-1B call
        that fell just above vs.\ just below the identity cutoff between DB
        vintages. Not investigated at the per-hit-identity level in the
        current rerun.
  \item \textbf{PR246B1C ST call.} \texttt{-} rather than ST2144 --- an
        allele-cutoff artefact, not evidence of a different lineage, but
        \emph{formally} the ST is not reproduced for that isolate.
\end{enumerate}

\subsection*{Threats to the replication result}
\begin{itemize}[leftmargin=1.4em]
  \item \emph{Assembly identity ambiguity.} We trust the authors' WGS
        accessions in PRJNA548686 / PRJNA412434 to correspond to the
        Table~1 isolate labels. No independent SNP-fingerprint check of
        assembly $\leftrightarrow$ isolate mapping was done.
  \item \emph{Database vintage confound.} Every observed discrepancy is
        consistent with 2021 vs.\ current ResFinder/NCBI panels. We did not
        rerun with a pinned 2021 DB snapshot, so we cannot separate
        ``paper called this wrongly'' from ``DB has since split/renamed
        this allele.''
  \item \emph{No plasmid-context reconstruction.} The paper's mobilome
        story (which plasmid Inc-type carries which $\beta$-lactamase) is
        gene-presence-replicated but not gene-\emph{location}-replicated;
        \texttt{abricate}+PlasmidFinder tell us Inc-types are present but
        not which contig they co-localise with which resistance gene.
        A hybrid-assembly / MOB-suite / long-read verification would be
        needed to fully replicate the plasmid-linkage claims, and is out of
        scope here.
  \item \emph{No phylogeny.} Author phylogroup calls are trusted from
        Table~1; we did not independently rerun ClermonTyping or a
        core-genome tree.
\end{itemize}

\subsection*{What a stronger replication would add}
\begin{enumerate}[leftmargin=1.4em]
  \item Pinned 2021 ResFinder DB snapshot to isolate DB-drift as the cause
        of CTX-M-15$\rightarrow$-55 and $\pm$1 co-carriage.
  \item Independent phylogroup call (ClermonTyping) and cgMLST.
  \item Long-read or Unicycler hybrid re-assembly of PN256E8 to resolve the
        TEM allele complement and confirm the CTX-M allele on-contig.
  \item Plasmid reconstruction (MOB-suite, plasmidSPAdes) to physically
        link CTX-M-15 to its Inc-type replicon --- the paper's central
        mobilome claim.
  \item Public-tree contextualisation (Enterobase HierCC / PathogenWatch)
        to test the pig-to-human clonal-transfer hypothesis raised by the
        original authors.
\end{enumerate}

\subsection*{Bottom line}
The paper's quantitative headline and its underlying isolate-level
attributions reproduce cleanly on the authors' deposited data with modern
open-source tooling. The residual discrepancies are (i) a single per-isolate
CTX-M allele call, (ii) a $\pm$1 co-carriage count, and (iii) one
MLST-identity edge case --- \emph{none} of which alter the paper's biological
conclusion. We therefore classify this as a genuine \textbf{PARTIAL
replication, borderline REPLICATED}, and flag DB-vintage control as the
single most impactful next step.

\end{document}
